\documentclass[addpoints,twoside]{exam}
\usepackage{problemset}

\begin{document}

\FontEasyRead

\Heading
{Microeconomics}
{Consumer Behavior \& Demand - Version B}
{\textbf{Instructions: }Answer each question clearly. Show reasoning where appropriate.}

\begin{questions}

    \Problem{Explain what it means for consumer preferences to be complete and transitive. Why are these assumptions important?}

    \Problem{Why can indifference curves not cross? What assumption would be violated if they did?}

    \Problem{Explain the concept of marginal rate of substitution (MRS). Why does MRS typically diminish as you move along an indifference curve?}

    \Problem{A consumer chooses between goods $X$ and $Y$.
        \begin{itemize}
            \item What happens to the budget line if the price of $X$ increases?
            \item How does the slope change?
        \end{itemize}}

    \Problem{Define marginal utility. Give an example of diminishing marginal utility in everyday life.}

    \Problem{Explain the condition
        \[
            \frac{MU_X}{P_X} = \frac{MU_Y}{P_Y}.
        \]
        Why must this hold at the optimal consumption bundle?}

    \Problem{A consumer has:
        \[
            MU_X = 18,\quad P_X = 3,\quad MU_Y = 12,\quad P_Y = 4
        \]
        \begin{itemize}
            \item Is the consumer maximizing utility?
            \item If not, how should they adjust consumption?
        \end{itemize}}

    \Problem{Suppose the marginal rate of substitution is 4.
        \begin{itemize}
            \item If $\frac{P_X}{P_Y} = 2$, is the consumer at an optimal bundle?
            \item Explain your reasoning.
        \end{itemize}}

    \Problem{Explain the substitution effect and income effect of a price decrease. Which effect always moves consumption in the same direction?}

    \Problem{Define a normal good and an inferior good. How would each appear on an income-consumption curve?}

    \Problem{What does an Engel curve show? How would it differ for a normal versus an inferior good?}

    \Problem{Explain how changes in price generate points on an individual demand curve.}

    \Problem{Suppose Consumer A demands 6 units at a given price and Consumer B demands 4 units.
        \begin{itemize}
            \item What is total market demand?
            \item What does “horizontal summation” mean in this context?
        \end{itemize}}

    \Problem{Why is the price elasticity of demand negative? Relate your answer to the law of demand.}

    \Problem{Explain why a linear demand curve does not have constant elasticity.}

    \Problem{A price decrease results in lower consumption of a good.
        \begin{itemize}
            \item What special case does this describe?
            \item What must be true about the income effect?
        \end{itemize}}

    \Problem{Explain consumer surplus in words. Why is it considered a measure of welfare?}

\end{questions}

\end{document}